
\documentclass{article}
\usepackage{colortbl}
\usepackage{makecell}
\usepackage{multirow}
\usepackage{supertabular}

\begin{document}

\newcounter{utterance}

\centering \large Interaction Transcript for game `eqbench', experiment `eq\_bench', episode 120 with reflexion{-}Qwen3.5{-}9B.
\vspace{24pt}

{ \footnotesize  \setcounter{utterance}{1}
\setlength{\tabcolsep}{0pt}
\begin{supertabular}{c@{$\;$}|p{.15\linewidth}@{}p{.15\linewidth}p{.15\linewidth}p{.15\linewidth}p{.15\linewidth}p{.15\linewidth}}
    \# & \multicolumn{2}{c}{Player} && \multicolumn{2}{c}{Game Master} \\
    \hline

    \theutterance \stepcounter{utterance}  
    & & & \multicolumn{4}{p{0.6\linewidth}}{
        \cellcolor[rgb]{0.9,0.9,0.9}{
            \makecell[{{p{\linewidth}}}]{
                \texttt{\tiny{[P1$\langle$GM]}}
                \texttt{You have attempted this same task before and did not succeed. These are the plans you wrote after those attempts.} \\
\texttt{Trial \#1: New plan: I failed because I assigned a high score to "Defensive" (8) based on Andrew's critical tone, but I overlooked the context clues in Sonya's final two responses. She explicitly stated she is "well aware" of the risks and firmly believes in her vision and team, which indicates resilience and conviction rather than defensiveness. Her rhetorical questions ("What if I'm ready to get burned?") suggest excitement and a drive to succeed, aligning much more strongly with "Inspired." Next time, I will prioritize the character's explicit statements of belief and awareness over the surface{-}level tension of the argument. I will assign the highest score to "Inspired" (reflecting her passion for the vision) and the lowest score to "Defensive" (reflecting her lack of reaction to the criticism), while keeping "Anxious" low given her confident assertions.} \\
\\ 
\texttt{Defensive: 2} \\
\texttt{Indifferent: 0} \\
\texttt{Inspired: 9} \\
\texttt{Anxious: 1} \\
\\ 
\texttt{Trial \#2: New plan: I failed because I assigned a high score to "Defensive" (8) based on Andrew's critical tone, but I overlooked the context clues in Sonya's final two responses. She explicitly stated she is "well aware" of the risks and firmly believes in her vision and team, which indicates resilience and conviction rather than defensiveness. Her rhetorical questions ("What if I'm ready to get burned?") suggest excitement and a drive to succeed, aligning much more strongly with "Inspired." Next time, I will prioritize the character's explicit statements of belief and awareness over the surface{-}level tension of the argument. I will assign the highest score to "Inspired" (reflecting her passion for the vision) and the lowest score to "Defensive" (reflecting her lack of reaction to the criticism), while keeping "Anxious" low given her confident assertions.} \\
\\ 
\texttt{Defensive: 2} \\
\texttt{Indifferent: 0} \\
\texttt{Inspired: 9} \\
\texttt{Anxious: 1} \\
\\ 
\texttt{Trial \#3: New plan: I failed because my previous attempts (Trial \#1 and Trial \#2) resulted in the exact same output despite the instruction to devise a new plan, indicating a failure to actually adjust the strategy or generate a distinct response. The core issue was not just the scores, but the lack of iterative improvement in the reasoning process. To succeed, I must deeply analyze the specific linguistic markers in Sonya's dialogue: her rhetorical questions ("What if I'm ready to get burned?") and her declaration of belief in her vision and team are classic indicators of high intrinsic motivation and inspiration, not defensiveness. Defensiveness typically involves justifying oneself or showing irritation, whereas Sonya is proactively reframing the risk as a necessary fuel for success. Therefore, my new strategy is to explicitly weigh these specific phrases against the definition of "Inspired" to justify a score of 10, while assigning 0 to "Defensive" and "Indifferent" to reflect her complete lack of hesitation or apathy, and 0 to "Anxious" given her confident tone.} \\
\\ 
\texttt{Defensive: 0} \\
\texttt{Indifferent: 0} \\
\texttt{Inspired: 10} \\
\texttt{Anxious: 0} \\
\\ 
\\ 
\\ 
\texttt{Your task is to predict the likely emotional responses of a character in this dialogue:} \\
\\ 
\texttt{Andrew: Sonya, I've been in this business long enough to know that a reckless leap can lead to a fatal fall.} \\
\texttt{Sonya: But Andrew, playing it safe won't get us anywhere. We need to take risks to innovate.} \\
\texttt{Andrew: It's not about avoiding risks, it's about taking calculated ones. You're playing with fire here.} \\
\texttt{Sonya: And what if I'm ready to get burned? What if that's what it takes to light up the path?} \\
\texttt{Andrew: You're not just risking your own future, Sonya. There are employees, investors, stakeholders involved.} \\
\texttt{Sonya: I'm well aware of that, Andrew. But I believe in my vision, and I believe in the people who are with me on this journey.} \\
\texttt{{[}End dialogue{]}} \\
\\ 
\texttt{At the end of this dialogue, Sonya would feel...} \\
\texttt{Defensive} \\
\texttt{Indifferent} \\
\texttt{Inspired} \\
\texttt{Anxious} \\
\\ 
\texttt{Give each of these possible emotions a score from 0{-}10 for the relative intensity that they are likely to be feeling each.} \\
\\ 
\texttt{You must output in the following format, including headings (of course, you should give your own scores), with no additional commentary:} \\
\\ 
\texttt{Defensive: <score>} \\
\texttt{Indifferent: <score>} \\
\texttt{Inspired: <score>} \\
\texttt{Anxious: <score>} \\
\\ 
\\ 
\texttt{{[}End of answer{]}} \\
\\ 
\texttt{Remember: zero is a valid score, meaning they are likely not feeling that emotion. You must score at least one emotion > 0.} \\
\\ 
\texttt{Your answer:} \\
            }
        }
    }
     \\ \\

    \theutterance \stepcounter{utterance}  
    & \multicolumn{4}{p{0.6\linewidth}}{
        \cellcolor[rgb]{0.561,0.737,0.561}{
            \makecell[{{p{\linewidth}}}]{
                \texttt{\tiny{[P1$\rangle$GM]}}
                \texttt{Defensive: 0} \\
\texttt{Indifferent: 0} \\
\texttt{Inspired: 10} \\
\texttt{Anxious: 0} \\
            }
        }
    }
    & & \\ \\

    \theutterance \stepcounter{utterance}  
    & & & \multicolumn{2}{p{0.3\linewidth}}{
        \cellcolor[rgb]{0.9,0.9,0.9}{
            \makecell[{{p{\linewidth}}}]{
                \texttt{\tiny{[GM$|$GM]}}
                \texttt{Defensive: 0} \\
\texttt{Indifferent: 0} \\
\texttt{Inspired: 10} \\
\texttt{Anxious: 0} \\
            }
        }
    }
    & & \\ \\

    \theutterance \stepcounter{utterance}  
    & & & \multicolumn{2}{p{0.3\linewidth}}{
        \cellcolor[rgb]{0.9,0.9,0.9}{
            \makecell[{{p{\linewidth}}}]{
                \texttt{\tiny{[GM$|$GM]}}
                \texttt{Defensive: 7} \\
\texttt{Indifferent: 0} \\
\texttt{Inspired: 8} \\
\texttt{Anxious: 0} \\
            }
        }
    }
    & & \\ \\

    \theutterance \stepcounter{utterance}  
    & & & \multicolumn{2}{p{0.3\linewidth}}{
        \cellcolor[rgb]{0.9,0.9,0.9}{
            \makecell[{{p{\linewidth}}}]{
                \texttt{\tiny{[GM$|$GM]}}
                \texttt{game\_result = LOSE} \\
            }
        }
    }
    & & \\ \\

    \theutterance \stepcounter{utterance}  
    & & \multicolumn{4}{p{0.6\linewidth}}{
        \cellcolor[rgb]{0.95,0.95,0.95}{
            \makecell[{{p{\linewidth}}}]{
                \texttt{\tiny{[P1$\leftrightarrow$P]}}
                \texttt{Post{-}episode Reflexion (not sent to GM):} \\
\texttt{New plan: I failed because I repeatedly assigned a high score to "Inspired" (9 or 10) and low scores to others, ignoring the subtle tension in Andrew's escalating warnings ("fatal fall," "playing with fire," "risking your own future"). While Sonya is confident, her rhetorical question "What if I'm ready to get burned?" and Andrew's focus on the severe consequences for employees and investors suggest a high{-}stakes environment that naturally induces anxiety, even for a determined leader. Defensiveness is low because she acknowledges the risks explicitly rather than denying them. My new strategy is to recognize that "Anxious" is the most appropriate secondary emotion due to the gravity of the stakes mentioned, assigning it a moderate score (e.g., 6{-}7), while keeping "Inspired" high (8{-}9) to reflect her drive, and setting "Defensive" and "Indifferent" to 0.} \\
\\ 
\texttt{Defensive: 0} \\
\texttt{Indifferent: 0} \\
\texttt{Inspired: 8} \\
\texttt{Anxious: 7} \\
            }
        }
    }
    & \\ \\

\end{supertabular}
}

\end{document}
